% arara: xelatex
%: {options: -output-directory=./output}
% Hinweis: Optionen der Dokumentenklasse werden an alle folgenden \usepackage{package} Befehle weitergegeben
\documentclass[
	fontsize=12pt,
	paper=a4,
	parskip=half,
	twoside=false,
	numbers=noenddot,	% Kein Punkt am Ende einer Überschrift
	%draft=true,			% Deckt Schwächen auf: overfull und full boxes werden markiert; Bilder werden nicht geladen
	bibliography=totoc,	% Literaturverzeichnis ins Inhaltsverzeichnis aufnehmen
	listof=totoc,		% Tabellen- und Abbildungsverzeichnis ins Inhaltsverzeichnis aufnehmen
	titlepage=true,		% Separate Titelseite; Gestaltung mit Hilfe der Titlepage-Umgebung
	headsepline=true,	% Kopflinie aktivieren
	footsepline=true,	% Fußlinie aktivieren
	abstracton			% Abstract aktivieren
]{scrreprt}

% Zeichenkodierung Input ist UTF-8: Umlaute können direkt eingegeben werden
\usepackage[utf8]{inputenc}

% Zeichenkodierung Ausgabe ist T1-Kodierung: Wichtig für die Ausgabe von Umlauten
\usepackage[T1]{fontenc}

% Schrift festlegen
\usepackage{fontspec}

\setmainfont[
BoldFont=Arial Bold.ttf,
ItalicFont=Arial Italic.ttf,
BoldItalicFont=Arial Bold Italic.ttf
]{Arial.ttf}

\usepackage{titlesec}
\renewcommand{\sectfont}{\rmfamily}
% Sprachauswahl für Lokalisierungen und Silbentrennung
\usepackage[german]{babel}

% Zitate: Anführungszeichen automatisch anhand der Sprache wählen
\usepackage[babel=true]{csquotes}

% Source-Code-Listings
\usepackage{listings}
\lstset{
    language=SQL,
    basicstyle=\ttfamily
} 

% BibTeX-Symbol
\usepackage{texnames}

%Farbpaket laden
\usepackage{xcolor}

% Symbole, z.B. Haken
\usepackage{pifont}

% Zeilen in Tabellen zusammenfassen
\usepackage{multirow}

% Silbentrennung kann bei bestimmten Wörten mit Hilfe von diesem Paket deaktiviert werden 
\usepackage{hyphenat}

% Abkürzungsverzeichnis
\usepackage[printonlyused, withpage]{acronym}
% Abstand mit Punkten füllen
\renewcommand*\bflabel[1]{\textbf{\normalsize{#1}}\hfill}

% Tiefe des Inhaltsverzeichnisses
\setcounter{tocdepth}{2}

% Punkte im Inhaltsverzeichnis
\usepackage{tocstyle}
\usetocstyle{allwithdot}

% Zum Einbinden von PDF-Dateien.
\usepackage{pdfpages}

% Paket zum Anpassen von Kopf- und Fußzeilen
\usepackage[plainfootsepline, plainheadsepline, headsepline, footsepline, automark]{scrpage2}
% Liniendicke
\setheadsepline{0.1pt}
\setfootsepline{0.1pt}

% Kopf- und Fusszeile löschen
\clearscrheadfoot
% Kopf- und Fusszeile aktivieren
\pagestyle{scrheadings}

% Kopf links
\ihead[\titel]{\titel}

% Fuss links
\ifoot[\verfasser]{\verfasser}
% Fuss rechts
\ofoot[\pagemark]{\pagemark}

% Grafiken einbinden
\usepackage{graphicx}
% Pfad zu den Grafiken
\graphicspath{{imagery/}}

% Seitenränder setzen
\usepackage[left=3.5cm, right=2.5cm, top=2.5cm, bottom=3cm]{geometry}
\renewcommand*\chapterheadstartvskip{\vspace*{0cm}}

% Zeilenabstand auf 1.5 setzen
\usepackage{setspace}
\onehalfspacing

% Literaturverzeichnis
%\usepackage[backend=bibtex,style=authoryear]{biblatex}
\usepackage[backend=bibtex,style=alphabetic]{biblatex}
\bibliography{literature/Literature.bib}
\renewcommand{\bibname}{Literaturverzeichnis}
%Referenz zu URLs
\usepackage{url}

% Glossar
\usepackage[toc,nonumberlist]{glossaries}

% Titel als Referenzierung verwenden
\usepackage{titleref}

% Währungen
\usepackage{textcomp}

% Fussnoten fortlaufend nummerieren.
\usepackage{chngcntr}
\counterwithout{footnote}{chapter}

% Persönliche Daten
\newcommand{\titel}{Was ist Web 2.0?}
\newcommand{\art}{Übung zum wissenschaftlichen Arbeiten}
\newcommand{\studienbereich}{Wirtschaft}
\newcommand{\studiengang}{Onlinemedien}
\newcommand{\verfasser}{Jill Zander}
\newcommand{\kurs}{ON13}
\newcommand{\ausbildungsbetrieb}{HITS gGmbH}
\newcommand{\abgabedatum}{02.03.2014}
\newcommand{\unterschrift}{\rule{5cm}{0.2pt}}

% Links- und PDF-Einstellungen
\usepackage[hidelinks]{hyperref}
\hypersetup{
	pdfauthor = {\verfasser},
	pdftitle = {\titel},
	pdfsubject = {\art},
	pdfkeywords = {},
	pdfstartview = {Fit},
	colorlinks = {false},
	breaklinks = {true},
	bookmarksopen = {true}
}

% Verhinderung von Schusterjunge und Hurenkind
\clubpenalty = 10000
\widowpenalty = 10000
\displaywidowpenalty = 10000

% Seitenzäler für große, römische Zahlen
\newcounter{RomanPagenumber}

% Abkürzungen
\newcommand{\dash}{d.\,h.}
\newcommand{\zB}{z.\,B.}


\begin{document}
	% Title page
	\thispagestyle{plain}

\begin{titlepage}
	\enlargethispage{4.0cm}

	\begin{center}		
		\raisebox{-.5\height}{\includegraphics[width=3cm]{imagery/visual4}}
		\hfill
		\raisebox{-.5\height}{\includegraphics[width=6cm]{imagery/dhbw_mosbach}}
		\\
		\vspace{3cm}
		\large{\textbf{Fachbereich: \studienbereich}}\\
		\vspace{0.5cm}
		\large{\textbf{Studiengang: \studiengang}}\\
		\vspace{1.0cm}
		\Large{\textsc{\textbf{\titel}}}\\
		\vspace{1.0cm}
		\large{\textbf{\art}}\\
		\vspace{1cm}
		\vspace{1cm}
		\vspace{3cm}
		\begin{tabular}{rl}
			Autor:					& \verfasser\\
			Kurs: 					& \kurs\\ 
			Ausbildungsbetrieb:		& \ausbildungsbetrieb\\ 
            Wiss. Ansprechpartner:  & \betreuer\\
			Abgabedatum:		& \abgabedatum\\
			\vspace{0.5cm}\\
			Unterschrift:				& \unterschrift\\
		\end{tabular} 
	\end{center}
\end{titlepage}

	
	% Ab hier große, römische Seitenzahlen
	\pagenumbering{Roman}
	
	%falls nicht der Titel der Arbeit oben stehen soll, bitte einkommentieren
	%\ihead[\headmark]{\headmark}
	
	% Sperrvermerk
	%\addchap{Restriction Note}
Diese Projektarbeit ist mit einem Sperrvermerk versehen und enthält vertrauliche Daten der \ausbildungsbetrieb. Eine Veröffentlichung oder Vervielfältigung dieser Arbeit, als Ganzes oder auch auszugsweise, ist ohne ausdrückliche Genehmigung der \ausbildungsbetrieb nicht zulässig. Die Einsicht in diese Arbeit ist unbefugten Personen und Organisationen nicht gestattet. Diese Arbeit darf nur den Korrektoren, der Studiengangsleitung und dem Prüfungsausschuss der DHBW Karlsruhe zugänglich gemacht werden. Der Sperrvermerk ist zeitlich befristet und erlischt frühestens 3 Jahre nach Abschluss der Projektarbeit bzw. 3 Jahre nach Abschluss des Prüfungsverfahrens zu dieser Projektarbeit.
	
	% Abstract
	%\include{content/Abstract}
	
	% Eidesstattliche Erklärung
	\addchap{Ehrenwörtliche Erklärung}
Ich erkläre hiermit eidesstattlich, dass ich die vorliegende Arbeit selbstständig und ohne Benutzung anderer als der angegebenen Hilfsmittel angefertigt habe.
Aus den benutzten Quellen direkt oder indirekt übernommene Gedanken habe ich als solche kenntlich gemacht.

Diese Arbeit wurde bisher in gleicher oder ähnlicher Form oder auszugsweise noch keiner anderen Prüfungsbehörde vorgelegt und auch nicht veröffentlicht.
\\
\vspace{2cm}
\\
\unterschrift\\
\noindent\verfasser
\\
\noindent Mosbach, den \today
	
	% Contentssverzeichnis
	\tableofcontents
	\pagebreak
	
	% Abbildungsverzeichnis
	%\listoffigures
	%\pagebreak
    
    % Code Listing Verzeichnis
    %\lstlistoflistings
    %\pagebreak
	
	% Tabellenverzeichnis
	%\listoftables
	%\pagebreak
	
	% Abkürzungsverzeichnis
	%\addchap{List of Abbreviations}
	%% Alle Abkürzungen hier definieren.
% Alphabetische Sortierung ist von Hand vorzunehmen.
\begin{acronym}[TT] % Längstes Akronym in eckigen Klammern
	% Kein Abstand, kompakte Darstellung
	%\setlength{\itemsep}{-\parsep}
	\acro{API}{Application Programming Interface}
    \acro{BAI}{Business Application Integration}
    \acro{EBM}{EleGant Business Model}
    \acro{MVC}{Model-View-Controller}
    \acro{P5}{Paradigm 5}
    \acro{RID}{Readable ID}
    \acro{RUID}{Readable Unique ID}
    \acro{SAP HANA XS}{SAP HANA Extended Application Services}
    \acro{SQL}{Structured Query Language}
	%\acrodefplural{API}[APIs]{Application Programming Interfaces}
\end{acronym}
	%\pagebreak
	
	% Abschnitt beenden
	\clearpage
	
	% Seitenzähler erhält den Wert der aktuellen groß römisch nummerierten Seite; Zwischenspeichern für später
	\setcounter{RomanPagenumber}{\value{page}}
	
	% Ab hier arabische Seitenzahlen
	\pagenumbering{arabic}
		
	% Kopf links auf aktuelles Kapitel ändern
	\ihead[\headmark]{\headmark}

% 1. Introduction
%   1.1 Motivation
%   1.2 Definition
%     a. Template
%     b. Engine
%     c. Template Engine
% 2. Demand
%   2.1 Usage Scenarios
%     a. XSJS Functions
%     b. Stored Procedures
%     c. Database Triggers
% 3. Architecture & Environment
%   3.1 Template
%     a. Token Principle
%     b. Meta Code
%   3.2 Template Engine
%     a. HANA XS
%     b. Generic Architecture
%     c. Required Classes & Entities
% (4. Implementation)
% 4. Conclusion

	
	\chapter{Einleitung}
Automatisieren ist immer günstiger
\section{Zielgruppe}
UX Designer, Softwareentwickler
Grundwissen; CRM, Softwareentwicklung, Datenbanken, 
\section{Problemstellung}
Wachsende Komplexität
Einfachere Interfaces
Eingabeformulare als Barriere
CRM-Software möglicherweise Hinterher
Efforts von Salesforce

\section{Motivation}
human indexing is expensive, cost reduction because expert-labour \cite{shneiderman}
non-tech affinity 
bei großen system fällt Überblick schwer.

	\chapter{Konzeption}
Konzept extraktion und separate speicherung als Datensatz
technologische Grundlagen beachten
alchemy api
learning classifiers
regular expressions
scope der arbeit beachten
größe des Trainingdatensatzes beachten
analogie zu theoretischer herangehensweise
problem des trainingsdatensatz: bereits exportierte, verfremdete daten
\chapter{Implementierung}
\section{Klassifikation der Eingabe}
\section{Extraktion von Information}

\chapter{Schluss}
learning by doing encouraged. people get introduced to the problem domain throught the suggestions served by the system.




	\clearpage
	
	% Ab hier weiter mit großen, römischen Seitenzahlen
	\pagenumbering{Roman}
	
	% Seitenzähler auf zwischengespeicherten Wert für große, römische Zahlen setzen
	\setcounter{page}{\theRomanPagenumber}

	% Kopf links auf Titel ändern
	\ihead[\titel]{\titel}	
	
	% Anhang
	% \addchap{Appendices}
	% \appendix
	% \input{content/Anhang}
	
	% % Beigaben
	% \addchap{Beigaben}
\label{cha:Beigaben}

Im Folgenden ist die Verzeichnisstruktur der beigelegten CD dargestellt.
 
\begin{description}
	\item[Abbildungen]
		Beinhaltet die verwendeten Abbildungen.
	\item[Bachelorarbeit]
		Beinhaltet die Bachelorarbeit und die Interviewprotokolle als \textit{PDF}.
		Der zugehörige \emph{\LaTeX{}} Source Code befindet sich im Unterverzeichnis \textit{LaTeX}.
		Im Unterverzeichnis \textit{DOC} befindet sich die von \textit{PDF} in \textit{DOC} konvertierte Version dieser Arbeit.
	\item[Literatur]
		Beinhaltet das Literaturverzeichnis im \textit{\BibTeX{}} Format, sowie die digital verfügbaren Quellen als \textit{PDF} oder \textit{MHT}.
\end{description}
	
	% % Glossar, falls benötigt
	% \newacronym{js}{JS}{JavaScript}
\newacronym{ie}{IE}{Informationsextraktion}
\newacronym{ae}{AE}{automatische Entitätserkennung}
\newacronym{crm}{CRM}{Costumer Relationship Management}
\newacronym{b2c}{B2C}{Business to Customer}
\newacronym{b2b}{B2B}{Business to Business}

% Entität
% Objekt
% Token


	% \makeglossaries
	% \printglossary
	
	% Literaturverzeichnis
	\printbibliography
	
	% Aktuellen Abschnitt beenden
	\clearpage
\end{document}
